% Generated draft: requires mathematical and source review.
\documentclass[11pt]{article}
\usepackage[margin=25mm]{geometry}
\usepackage{fontspec}
\setmainfont{Latin Modern Roman}
\usepackage{amsmath,amssymb,mathtools,unicode-math}
\setmathfont{Latin Modern Math}
\usepackage{enumitem,needspace,xurl,hyperref}
\setlength{\parindent}{0pt}
\setlength{\parskip}{5pt}
\setlength{\emergencystretch}{3em}
\newcommand{\N}{\mathbb{N}}
\newcommand{\Z}{\mathbb{Z}}
\newcommand{\Q}{\mathbb{Q}}
\newcommand{\R}{\mathbb{R}}
\newcommand{\C}{\mathbb{C}}
\newcommand{\F}{\mathbb{F}}
\newcommand{\A}{\mathbb{A}}
\newcommand{\PP}{\mathbb{P}}
\newcommand{\E}{\mathbb{E}}
\newcommand{\eps}{\varepsilon}
\newcommand{\dd}{\,\mathrm{d}}
\begin{document}
\section*{UnsolvedMath ID: 701329}
\textbf{Source problem code:} catalogue entry
\subsection{Definitions and mathematical statement}
Under the usual convention that an ordered group has a total order invariant under both left and right multiplication, let \((G_i,\leq_i)\), for \(i\in\{1,2\}\), be ordered groups. Thus, for all \(a,b,x,y\in G_i\),
\[
a\leq_i b\implies xay\leq_i xby.
\]
Let \(H_i\leq G_i\) carry the restricted order, and let \(\varphi:H_1\to H_2\) be a group isomorphism that is also an order isomorphism:
\[
h\leq_1 k\quad\Longleftrightarrow\quad\varphi(h)\leq_2\varphi(k)\qquad(h,k\in H_1).
\]
Form the free product with amalgamation
\[
P=G_1*_{\varphi}G_2
=(G_1*G_2)\big/\big\langle\!\big\langle h\varphi(h)^{-1}:h\in H_1\big\rangle\!\big\rangle,
\]
where the displayed elements are understood through the canonical inclusions into the free product and the double brackets denote normal closure.

A lattice order on \(P\) is a partial order \(\preccurlyeq\) such that every pair of elements has a greatest lower bound and a least upper bound, and
\[
a\preccurlyeq b\implies xay\preccurlyeq xby\qquad(a,b,x,y\in P).
\]
The group \(P\) is lattice orderable if it admits such an order. It is orderable, in the convention used here, if it admits a total order invariant under both left and right multiplication.

\textbf{Question (Bludov--Glass).} Do there exist \((G_1,\leq_1)\), \((G_2,\leq_2)\), \(H_1\), \(H_2\), and \(\varphi\) as above such that \(P\) admits a lattice order but admits no total order invariant under both left and right multiplication?

\textbf{Unresolved qualifications.} The supplied statement does not explicitly specify whether ``ordered'' means invariance under both left and right multiplication. The formulation above uses that standard convention. It also interprets orderability and lattice orderability as properties of the abstract amalgamated product; the supplied statement does not explicitly say whether an order on \(P\) must extend the prescribed orders on its factors. These conventions require mathematical review.
\subsection{Short English statement}
Can two totally ordered groups, amalgamated along subgroups identified by an order-preserving group isomorphism, yield a group admitting a translation-invariant lattice order but no translation-invariant total order? This formulation uses orders invariant under both left and right multiplication; the intended convention and any requirement to extend the factor orders need confirmation.
\subsection{Sources}
\begin{itemize}
\item Local UnsolvedMath registry, supplied complete source record: permanent ID \(701329\), exact source problem code \texttt{catalogue entry}, title \emph{Kourovka Notebook Problem 17.27}.
\item V. V. Bludov and A. M. W. Glass, problem proposed in 2010, \emph{The Kourovka Notebook (Unsolved Problems in Group Theory)}, 17th issue (2010), Problem 17.27. The supplied record also references the 21st issue edition (2026), \texttt{arXiv:1401.0300v48}; the collection description names E. I. Khukhro and V. D. Mazurov as editors.
\item The supplied background and status review report an affirmative answer labelled ``Unverified AI claim,'' attributed to Achyuth Jayadevan, \emph{Preprint} (2026): \url{https://zenodo.org/records/22937468}. The status review says that this claim is marked with \(\#\) in the October 2026 version of the 21st issue. This is a supplied report of an unverified claim, not an established resolution.
\item Supplied status-review date: \texttt{2026-10-16}. The registry retains the label \texttt{open}, with evidence classified as \texttt{unrefereed\_claim}. This assessment is supplied dated curation and has not been independently verified.
\end{itemize}
\end{document}
